\chapter*{Appendices}
\addcontentsline{toc}{chapter}{Appendices}

\section*{Appendix A: Reproducibility and Execution Guide}
\addcontentsline{toc}{section}{Appendix A: Reproducibility and Execution Guide}

This appendix provides supporting information for reproducing the experimental workflow used in this thesis. The experiments were implemented as Jupyter notebooks and executed locally using the same document corpus, model configuration, scenario questions, logging format, and evaluation scripts.

The main notebooks used in the experiments were:
\begin{itemize}
    \item \texttt{rag\_basic.ipynb}
    \item \texttt{rag\_vectordb.ipynb}
    \item \texttt{rag\_vectordb\_tuned.ipynb}
    \item \texttt{rag\_reranker.ipynb}
    \item \texttt{rag\_hybrid.ipynb}
    \item \texttt{rag\_hybrid\_\&\_rerank.ipynb}
\end{itemize}

The local execution environment used Ollama for model serving and Milvus for the vector-database and hybrid-retrieval pipelines. The generation model was \texttt{llama3.1}, while the embedding model was \texttt{mxbai-embed-large}. Each notebook followed the same general workflow: loading the document corpus, splitting documents into chunks, indexing or storing embeddings, executing the scenario questions, generating answers, and saving results to the shared results file.

This appendix is intended to support reproducibility by documenting the practical execution order and implementation files used to generate the experimental outputs discussed in the main thesis.

\newpage
\section*{Appendix B: Evaluation Scenarios and Scoring}
\addcontentsline{toc}{section}{Appendix B: Evaluation Scenarios and Scoring}

This appendix lists the scenario questions used to evaluate the implemented RAG pipelines. The same questions were used across all pipeline variants to ensure that differences in retrieval quality, answer quality, and execution time were caused by the retrieval strategy rather than by changes in the input task.

\begin{enumerate}
    \item \textbf{What is the first-response SLA for a Sev-1 support incident?}
    
    This scenario tests direct factual retrieval from a single relevant document.

    \item \textbf{What steps must be completed before a new employee is considered ready for day one at OrchestrAI?}
    
    This scenario tests whether the system can retrieve and summarize a procedural checklist.

    \item \textbf{A new engineer starts on Monday, but their laptop has not arrived and their required access is still pending. Which teams are involved and what should happen next?}
    
    This scenario tests cross-document retrieval involving HR, IT, and access-related responsibilities.

    \item \textbf{Can OrchestrAI make an urgent hardware purchase without following the normal PO workflow if a warehouse outage is at risk?}
    
    This scenario tests policy retrieval and exception-based reasoning.

    \item \textbf{A customer renewal is at risk because support cases are recurring, a requested feature is still unavailable, and replacement hardware is delayed. What should OrchestrAI do next, and which teams must be involved?}
    
    This scenario tests multi-hop retrieval across support, product, logistics, and sales-related information.
\end{enumerate}

Manual answer quality was scored using three values: 1.0 for a correct answer, 0.5 for a partially correct answer, and 0.0 for an incorrect or unsupported answer. Retrieval quality was evaluated using whether the expected supporting documents appeared in the retrieved sources, together with precision, recall, and F1. Automated answer-quality metrics included ROUGE-L, BERTScore, sentence-embedding semantic similarity, and answer length.